\documentclass[11pt,a4paper]{article}
\usepackage[top=42pt,bottom=42pt,left=70pt,right=70pt]{geometry}  % to make pages wider

\usepackage{amsmath}
\usepackage{amsfonts}
\usepackage{graphicx}  % for figures etc.
\usepackage{url}  % to handle URLs 
\usepackage[comma,authoryear]{natbib}  % for author-date ref style --  allows \citet (textual), and \citep (parenthesized)
\usepackage[hidelinks]{hyperref}
\hypersetup{allcolors=[rgb]{0, 0, 0.33},colorlinks=true}

\renewcommand {\cite} {\citep}  % default for cite is citet in natbib - changes it to citep
\renewcommand\bibname{References}    % if not using newrucsthesis sty file

\usepackage{booktabs}
\usepackage{array}
\usepackage{enumitem}
\usepackage{listings}
\usepackage{amssymb}

\usepackage{setspace}
\doublespacing

% Optional: Customizing how the code looks
\lstset{
    basicstyle=\ttfamily\small,
    numbers=left,           
    numberstyle=\tiny,     
    stepnumber=2,          
    numbersep=5pt,
    frame=single,           
    rulecolor=\color{black},
    frameround=tttt,        
    breaklines=true         
}

\title{Results and Analysis}
\author{g23n7880 }
\date{September 2026}

\begin{document}

\maketitle
\section{Overview}

This chapter presents the results of three evaluations conducted on the semi-automated Sesotho corpus construction pipeline. The evaluations are organised around the three research questions and their associated success criteria. 

The first evaluation addresses the orthographic classifier's  performance while the second and third evaluations addresses the corpus's downstream utility.

\begin{table}[ht]
    \centering
    \begin{tabular}{llll}
     \toprule
    \textbf{RQ} & \textbf{Evaluation} & \textbf{Criterion} &
    \textbf{Outcome} \\
    \midrule
    RQ2 & Orthographic classifier & Macro F1 $\geq$ 0.85 & Met (0.87) \\
    RQ3 & Perplexity (FLORES-200) & Reduction, $p < 0.05$ &
          Met (31.8\%, $p < 0.001$) \\
    RQ3 & NER (NWU corpus) & F1 improvement, $p < 0.05$ &
          Not met ($\Delta = -0.001$) \\
    \bottomrule     
    \end{tabular}
    \caption{Mapping of evaluations to research questions and success criteria.}
    \label{tab:results_map}
\end{table}

\section{Orthographic classifier evaluation}

\subsection{Experimental Setup }

The orthographic classifier was evaluated on a stratified 20\% held-out test set of 1726 sentences (1073 SAS and 653 LS), drawn from the same verified corpus used for training. A rule-based baseline and trained character n-gram logistic regression classifier were evaluated on the same test set. The rule-based baseline implements the substitution rules documented by \citet{makutoane2022people} as a marker-counting function applied to pre-normalised text.

\subsection{Rule-based Baseline Results}

The rule-based baseline achieved a macro F1 of 0.4111 on the held-out test set. A per-class breakdown reveals that LS precision was 1.00 while LS recall was 0.03, meaning the baseline correctly identified almost no LS sentences despite perfect precision on the few it did label as LS. SAS precision was 0.63 and recall was not reported as the class dominated prediction by default.

This result is informative as it confirms the explicit substution rules documented are insufficient as a sentence-level classification mechanism. The markers such as \textit{ea}, \textit{oa}, \textit{ya}, \textit{wa} do not appear uniformly across all Sesotho sentences. Many grammatically complete sentences that rely exclusively on the presence of these markers will therefore fail to identify the majority of genuine LS sentences which  motivates the need for a learned classifier that can exploit distributional character patterns beyond the documented marker set.

\subsection{Learned classifier results}
The character n-gram logistic regression classifier achieved a macro F1 of 0.8714 on the held-out test set, exceeding the proposed threshold of 0.85 and surpassing the rule-based baseline by 0.46 F1 points.

\begin{table}[ht]
    \centering
    \begin{tabular}{lrrrrr}
     \toprule
    \textbf{Method} & \textbf{Prec LS} & \textbf{Rec LS}
                    & \textbf{Prec SAS} & \textbf{Rec SAS}
                    & \textbf{Macro F1} \\
    \midrule
    Rule-based (Makutoane) & 1.00 & 0.03 & 0.63 & n/a & 0.41 \\
    Char n-gram + LR       & 0.80 & 0.89 & 0.93 & 0.87 &
                             \textbf{0.87} \\
    \midrule
    Proposed target        &      &      &      &      & 0.85 \\
    \bottomrule     
    \end{tabular}
    \caption{Classification performance on 1726 sentence held-out test set.}
    \label{tab:clf_compare}
\end{table}

The confusion matrix shows 583 of 653 LS sentences correctly classified and 930 of 1073 SAS sentences correctly classified, with 70 LS sentences misclassified as SAS and 143 SAS sentences misclassfied as LS. The asymmetry in misclassification reflects the class imbalance in the training data, which the balanced class weighting partially but not fully corrects.

\subsection{Interpretability Analysis}

Inspection of the logistic regression coefficient vector revealed which character n-gram features most strongly predicted each orthographic variant. The top LS predicting feature was \texttt{oa} with a coefficient of $-3.44$ and the top SAS predicting features included \texttt{wa} with a coefficient of $+1.67$ which corresponds 



\end{document}
